\section{Summary}

Proposed compression algorithm shows promise in the lossless benchmarks, as on kodak dataset FRIF achieved overall better
compression rates than ubiquitous PNG format. Novel context modeling and prediction greatly enhanced the ability of the compressor to find spatial redundancies
and lower the amount of bits required to signal one pixel. This advantage quickly diminishes versus JPEG2000, a similar wavelet compressor, that
tacross all lossless benchmarks performed better. 
For lossy mode experiments showed that proposed quantization pattern produces pixelated versions of the original image, damaging the quality observed
via SSIM2 metric.

\subsection{Room for improvement}
There is a noticable room for improvement for presented codec and areas worth exploring to enhance both lossless and lossy compression ratios. 
One of them can be the choice of used wavelet in DWT, as Haar wavelet is the simplest wavelet that fitted the recursive tame-twindragon fractal nature. 
A wavelet more suited for image compression could benefit lossless mode, bringing FRIF closer to JPEG2000 performance.
The quantization scheme could potentially be enhanced by an approach similar to JPEGXL Modular mode with a tendency term, that limits the pixelation of dead-zoned
parameters. For a small compression ratio boost FRIF with better width prediction for context modelling could abandon saving \texttt{off\_distribution\_values} array as the image header.

\subsection{Closing remarks}
This thesis presented FRIF, a new image compression codec based on tame-twindragon fractal grid decomposition and discrete wavelet transform
with novel context modeling scheme. Although the lossy encoding mode requires additional refinement, FRIF gives a decent backbone of a codec that
can losslessly compress images with on-pair performance to PNG. Proposed solutions in grid selection, context modeling and prediction
with improvements mentioned in the thesis could yield even better results for both lossy and lossless compression, leaving space for further research. 
Presented compression scheme and implementation details can be found on an open-source repository hosted on GitHub under following URL: \url{https://github.com/pagmerek/frave}
